\pdfoutput=1
\documentclass[11pt,a4paper]{article}
\usepackage{abhijit-research}
\renewcommand{\PaperCode}{P00}
\renewcommand{\PaperKind}{RESEARCH PAPER}
\renewcommand{\PaperVersion}{VERSION 1.1}
\renewcommand{\PaperDate}{AUGUST 2026}
\renewcommand{\PaperTitle}{Recurrence Exponents}
\renewcommand{\PaperSubtitle}{Accessibility Growth and Predictive-State Promotion}
\renewcommand{\PaperFullTitle}{Recurrence Exponents, Accessibility Growth, and Predictive-State Promotion}
\renewcommand{\PaperShortTitle}{Recurrence Exponents}
\renewcommand{\PaperKeywords}{Lyapunov exponent; predictive state; representation repair; quantum channel; criticality; accessibility}
\renewcommand{\PaperClassification}{The recurrence theorems are exact under displayed cocycle assumptions. Cross-domain identifications require explicit intertwiners. Empirical sections establish finite future relevance rather than asymptotic exponents.}
\renewcommand{\PaperCitation}{Singh, Abhijit. 2026. \textit{Recurrence Exponents, Accessibility Growth, and Predictive-State Promotion}. Research manuscript, version 1.0.}
\begin{document}
\MakeResearchTitle
\MakeResearchContents
\MakeResearchAbstract{%
A residual left by an insufficient representation is assigned a cocycle and a typed growth exponent. Exact finite-dimensional theorems show that negative exponent gives exponential assimilation and positive exponent gives exponential escape. The zero-exponent class is proved to be a critical band requiring a secondary polynomial or persistence index. An independent accessibility exponent measures growth of source multiplicity hidden by the available continuation algebra. The calculus is realized in primitive Markov laws and quantum channels, related through explicit bridges to arithmetic capacity exponents, and compared with two finite held-out representation-restoration studies. Stationary witness margin, operator-ideal regularity, and finite empirical gain are separated from the recurrence exponent rather than conflated with it. Positive escape becomes an exact trigger for predictive-state enlargement. A new exact finite realization is included: a nine-state carrier with a 36-cell pair-relation frontier in which side-forgetting produces a fixed report while the complete source relation has critical period-two recurrence.
}
\section{Residual dynamics}
A reporter compresses source histories into a current state. When a newly admitted continuation separates two histories that the reporter merged, the difference is a residual of the present representation. Let $B_0$ be a residual and let $G_n$ be the linearized residual propagator between representation updates. With intrinsic increments $\Delta\tau_n>0$,
\begin{equation}
B_n=G_{0:n}B_0,\qquad G_{0:n}=G_{n-1}\cdots G_0,
\qquad \tau_n-\tau_0=\sum_{j<n}\Delta\tau_j.
\end{equation}
Define
\begin{equation}
\boxed{\chi(B_0)=\limsup_{n\to\infty}
\frac{\log\|G_{0:n}B_0\|-\log\|B_0\|}{\tau_n-\tau_0}.}
\end{equation}
The exponent is attached to a declared residual, propagator, norm class, and clock. It is not a free-floating label for any quantity that happens to have a boundary at one half.

\section{Finite-dimensional theorem}
In finite dimension all norms are equivalent, so $\chi$ is norm independent for a fixed clock.
\begin{theorem}[Constant-operator formula]
Let $G$ be a linear operator, $\tau_n=n$, and $V_B$ the cyclic invariant subspace generated by nonzero $B$. Then
\begin{equation}
\chi(B)=\log r(G|_{V_B}),
\end{equation}
where $r$ is spectral radius.
\end{theorem}
\begin{proof}
Jordan form expresses $G^nB$ as a finite sum of terms $n^k\lambda^n$. Polynomial factors vanish after logarithm and division by $n$; the largest modulus eigenvalue with nonzero component in $B$ remains.
\end{proof}

\begin{corollary}[Exponential trichotomy]
If $\chi(B)<0$, the residual is exponentially assimilated. If $\chi(B)>0$, it grows exponentially along a subsequence and eventually exceeds every bounded residual budget. The case $\chi(B)=0$ is the critical exponential band.
\end{corollary}

\section{The zero-exponent band}
Zero exponential rate is not one dynamical behaviour. Define
\begin{equation}
\beta(B)=\limsup_{n\to\infty}
\frac{\log(\|B_n\|/\|B_0\|)}{\log(1+n)},
\qquad
p_\infty(B)=\limsup_{n\to\infty}\|B_n\|.
\end{equation}
Within $\chi=0$:
\begin{itemize}
\item $\beta<0$ gives subexponential assimilation;
\item $\beta>0$ gives subexponential escape;
\item $p_\infty>0$ gives retained contrast;
\item bounded unitary or oscillatory modes have $\chi=0$ and nonzero persistence.
\end{itemize}
The selection law is therefore: assimilate when $\chi<0$; promote when $\chi>0$; and run a secondary critical-band test when $\chi=0$.

\section{Accessibility exponent}
Residual growth and inaccessible provenance are different. For a resource-bounded continuation family $\Lambda(B_\tau)$, define
\begin{equation}
\eta=\limsup_{\tau\to\infty}\frac1\tau
\log\mu\bigl([F]_{\Lambda(B_\tau)}\bigr),
\end{equation}
where the equivalence class contains source occurrences indistinguishable to the available continuation algebra. The pair $(\chi,\eta)$ distinguishes:
\begin{enumerate}
\item $\chi<0$, $\eta\approx0$: transparent assimilation;
\item $\chi<0$, $\eta>0$: local relaxation with growing hidden provenance;
\item $\chi\approx0$, $\eta>0$: critical recurrence with growing inaccessible correlation;
\item $\chi>0$: state-space failure and representation genesis.
\end{enumerate}
This two-coordinate system is required for horizon and thermodynamic interpretations in which geometry can be stable while accessibility changes.

\section{Primitive Markov realization}
For a primitive column-stochastic matrix $C$ with stationary vector $\pi$, let
\begin{equation}
\Pi=\pi\mathbf1^{\mathsf T},\qquad R=C-\Pi.
\end{equation}
Then $C^n=\Pi+R^n$ and every zero-sum contrast $d$ obeys
\begin{equation}
\chi(d)=\log r(R|_{V_d})<0.
\end{equation}
For the corrected four-state repeated instrument, the nontrivial eigenvalues are
\begin{equation}
0.35301358,\qquad0.27498875,\qquad0.19041902.
\end{equation}
The slowest contrast exponent is
\begin{equation}
\chi_{\rm mark}=\log(0.35301358)\approx-1.0412
\end{equation}
per use. This explains exponential loss of source memory while the stationary entropy rate remains positive.

\section{Quantum-channel realization}
For a primitive quantum channel $\mathcal T$, a traceless eigenoperator $X$ with eigenvalue $\lambda$ has
\begin{equation}
\chi(X)=\log|\lambda|.
\end{equation}
A product of unread projective channels contracts toward the common commutant. Under irreducibility only the identity remains fixed, so every traceless state contrast has $\chi<0$.

A contextual witness margin
\begin{equation}
\Delta_{\rm wit}=q-\beta_{\rm NC}
\end{equation}
is a different coordinate. A channel can have $\chi_{\rm state}<0$ while sustaining a positive state-independent witness. The complete protocol class is therefore at least two-dimensional:
\begin{equation}
(\chi_{\rm state},\Delta_{\rm wit}).
\end{equation}

\section{Arithmetic realizations}
The orientation operator has a Schatten boundary at $\Re s=1/2$, but this operator-regularity line is not itself a recurrence exponent. The recurrence coordinates are separately defined:
\begin{equation}
\chi_{\rm cap}=\limsup_{N\to\infty}\frac{\log(A_N/N)}{\log N},
\end{equation}
for the weighted capacity excess, and
\begin{equation}
\chi_{\rm RH}=\limsup_{N\to\infty}\frac1N\log(1+\Theta_N)
\end{equation}
for the causal-tail escape representation. Within the frozen arithmetic bridge, RH is equivalent to $\chi_{\rm cap}=0$ and to absence of positive transverse escape. The half-cost operator boundary and recurrence zero line coincide only after the explicit source-preserving bridge, not by visual similarity.

\section{Finite empirical promotion}
The relation-aware representation study and the chronological channel-restoration study prove future relevance on finite held-out data. They do not estimate an asymptotic exponent unless a repeated residual propagation law is declared.

Three independent representation rates $92.96\%$, $90.74\%$, and $91.48\%$ combine relationally to $97.78\%$. The omitted environmental channel changes untouched-future coordinate accuracy from $72.0836\%$ to $74.7125\%$ and exact next-vector accuracy from $3.7540\%$ to $7.2492\%$. These are finite promotion diagnostics: the discarded contrast was not future-inert.

\section{Selection theorem}
\begin{theorem}[Recurrence-based state selection]
Let a residual evolve under a fixed finite-dimensional cocycle, and suppose state promotion has a declared finite cost. Then:
\begin{enumerate}
\item $\chi<0$ is sufficient for exponential assimilation;
\item $\chi>0$ is sufficient for eventual failure of every bounded residual budget and therefore requires state promotion or rejection of the model;
\item $\chi=0$ requires a secondary persistence or polynomial-rate test before either assimilation or promotion is justified.
\end{enumerate}
\end{theorem}
The theorem is exact under the displayed assumptions. Cross-domain unification requires an explicit intertwiner identifying each domain's residual, propagator, clock, and norm class.

\section{Scientific claim}
The genuine result is not that one familiar Lyapunov exponent has been renamed. It is the placement of recurrence analysis inside a source-preserving representation calculus: residuals are generated by false closure, the exponent classifies whether the current representation can assimilate them, and positive escape becomes a mathematically explicit trigger for state-space enlargement. The cross-domain instances demonstrate the breadth of the calculus; they do not erase their local assumptions.


\section{Detailed recurrence calculus}
Let $V$ be a finite-dimensional normed space, let $G_n:V\to V$ be the residual propagator between two representation updates, and let the intrinsic clock satisfy $\Delta\tau_n>0$. The residual evolution is
\begin{equation}
B_n=G_{0:n}B_0,\qquad G_{0:n}=G_{n-1}\cdots G_0,
\end{equation}
and the recurrence exponent is
\begin{equation}
\chi(B_0)=\limsup_{n\to\infty}
\frac{\log\|G_{0:n}B_0\|-\log\|B_0\|}{\tau_n-\tau_0}.
\end{equation}
Because all norms are equivalent in finite dimension, changing the norm changes the numerator by a bounded additive term and therefore leaves $\chi$ invariant whenever $\tau_n\to\infty$.

\begin{theorem}[Constant-operator formula]
For $G_n=G$ and $\tau_n=n$, let $V_{B}$ be the cyclic invariant subspace generated by $B$. Then
\begin{equation}
\chi(B)=\log r(G|_{V_B}),
\end{equation}
where $r$ is spectral radius.
\end{theorem}
\begin{proof}
Put the restriction of $G$ to $V_B$ in Jordan form. Every component of $G^nB$ has the form $p(n)\lambda^n$, where $p$ is a polynomial whose degree is smaller than the corresponding Jordan-block size. After taking a logarithm and dividing by $n$, the polynomial term contributes $o(1)$. The eigenvalue of largest modulus that has nonzero projection from $B$ therefore determines the limit superior.
\end{proof}

\subsection{The zero-exponent band}
The sign trichotomy is exact at exponential scale, but $\chi=0$ contains several different behaviours. Define
\begin{equation}
\beta(B)=\limsup_{n\to\infty}
\frac{\log(\|B_n\|/\|B_0\|)}{\log(1+n)},
\qquad
p_\infty(B)=\limsup_{n\to\infty}\|B_n\|.
\end{equation}
Then polynomial decay has $\chi=0$ and $\beta<0$; bounded nonzero oscillation has $\chi=0$ and $p_\infty>0$; polynomial growth has $\chi=0$ and $\beta>0$. A scientifically usable selection rule is therefore
\begin{center}
\begin{tabular}{lll}
\toprule
condition&dynamical reading&representation action\\\midrule
$\chi<0$&exponential assimilation&compress after backward tests\\
$\chi>0$&exponential escape&promote or reject the state model\\
$\chi=0,\ \beta<0$&subexponential assimilation&retain until the decay certificate closes\\
$\chi=0,\ p_\infty>0$&critical retention&preserve as an active coordinate\\
$\chi=0,\ \beta>0$&subexponential escape&promote under the declared resource budget\\
\bottomrule
\end{tabular}
\end{center}

\section{Exact finite realizations}
\subsection{Primitive Markov law}
For a primitive column-stochastic matrix $C$ with stationary vector $\pi$, define
\begin{equation}
\Pi=\pi\mathbf1^{\mathsf T},\qquad R=C-\Pi.
\end{equation}
Then $\Pi^2=\Pi$, $\Pi R=R\Pi=0$, and
\begin{equation}
C^n=\Pi+R^n.
\end{equation}
Every zero-sum contrast evolves only through $R$. In the corrected four-record instrument, the nontrivial eigenvalues are
\begin{equation}
0.35301358,\qquad0.27498875,\qquad0.19041902,
\end{equation}
so the slowest exact exponential rate is
\begin{equation}
\chi_{\rm mark}=\log(0.35301358)\approx-1.0412
\end{equation}
per reuse. This explains exponential loss of source memory; it does not imply zero entropy production, because the stationary instrument continues to emit outcomes.

\subsection{Quantum channel and witness margin}
For a primitive quantum channel $\mathcal T$, a traceless eigenoperator $X$ with eigenvalue $\lambda$ has $\chi(X)=\log|\lambda|$. The stationary contextual-witness margin
\begin{equation}
\Delta_{\rm wit}=q-\beta_{\rm NC},\qquad q=\Tr(\sigma R_{\rm score}),
\end{equation}
is a second coordinate. A channel may satisfy $\chi_{\rm state}<0$ while $\Delta_{\rm wit}>0$, as in a state-independent contextual witness. Contraction, raw entropy, and certification are therefore not interchangeable names for one scalar.

\subsection{Arithmetic criticality}
The squarefree orientation operator has singular values $n^{-\sigma}$ and is Hilbert--Schmidt exactly for $\sigma>1/2$. This is an operator-ideal threshold. The capacity recurrence coordinate
\begin{equation}
\chi_{\rm cap}=\limsup_{N\to\infty}\frac{\log(A_N/N)}{\log N}
\end{equation}
asks a different question: whether the coercive excess grows subpower or by a positive power. The two boundaries are related by the complete arithmetic construction, but they are not identical merely because both place special structure at one half.

\section{Nine-state critical-return realization}
The prime-observer carrier gives an exact finite realization of critical retention. Write
\begin{equation}
A=\{1\}\sqcup\bigl(\{2,3,5,7\}\times\{+,-\}\bigr),
\end{equation}
where the four two-sided sectors correspond to $(2,9)$, $(3,8)$, $(5,6)$, and $(7,4)$. The return map flips the side and preserves the prime-sector label. For
\begin{equation}
r_q=e_{q,+}-e_{q,-},
\end{equation}
the return satisfies $Gr_q=-r_q$, so $\|G^nr_q\|=\|r_q\|$ and $\chi(r_q)=0$. A reporter retaining only $q$ sees a fixed point; the full state has period two. This is a strict example of
\begin{equation}
U(r_q)=0\quad\not\Rightarrow\quad\chi(r_q)<0.
\end{equation}
The complete pair-relation frontier of the nine states has $\binom92=36$ cells, decomposed as eight unit relations, four return relations, and twenty-four transverse relations. The finite theorem is developed in P17. Its role here is to distinguish report-level cancellation from dynamical assimilation.

\section{Finite promotion evidence and inference boundary}
The relation-aware representation experiment and the chronological feature-restoration experiment establish that a discarded difference changed an untouched future. That is a finite promotion certificate. It is not an estimate of an asymptotic $\chi$ unless a repeated residual propagation law is separately specified and measured. This distinction is essential: the exponent is a property of a dynamical cocycle, while held-out improvement is evidence that a particular quotient was insufficient.

\section{Falsification programme}
The recurrence-selection proposal fails as a cross-domain law if no source-preserving map sends the domain residual and its continuation to the displayed cocycle; if the sign is not invariant under admissible clock changes; or if a promoted residual repeatedly fails to alter later continuation after ablation. These are operational failure conditions rather than editorial reservations.

\begin{thebibliography}{99}
\bibitem{meyer} C. D. Meyer, Matrix Analysis and Applied Linear Algebra, SIAM, 2000.
\bibitem{seneta} E. Seneta, Non-negative Matrices and Markov Chains, Springer, 2006.
\bibitem{wolf} M. M. Wolf, Quantum channels and operations: a guided tour, lecture notes, 2012.
\bibitem{oseledets} V. I. Oseledets, A multiplicative ergodic theorem, Transactions of the Moscow Mathematical Society 19 (1968), 197--231.
\bibitem{crutchfield} J. P. Crutchfield and K. Young, Inferring statistical complexity, Physical Review Letters 63 (1989), 105--108.
\end{thebibliography}
\end{document}
